% VI21-CHALLENGE-SOLUTION: C01
\paragraph{Solution to Challenge VI.21.C1.}
\begin{solution}
For an absolute $K$-valued point, define
\[
\Phi:X(K)=\operatorname{Hom}(\Spec K,X)
\longrightarrow
\coprod_{x\in X}\operatorname{Hom}_{\mathrm{Fields}}(\kappa(x),K)
\]
by sending $u$ to its image point $x=u(*)$ together with the residue-field map
$\kappa(x)\to K$ induced from the local stalk map.  The field-valued point theorem proves that
this construction is bijective: given $(x,\iota)$, compose
\[
\Spec K\xrightarrow{\Spec\iota}\Spec\kappa(x)
\xrightarrow{i_x}X.
\]
The disjoint union is essential because the image point is part of the data.

Locality enters exactly when passing from $\OO_{X,x}\to K$ to $\kappa(x)\to K$.  Since the target
field has maximal ideal $(0)$, locality forces the inverse image of $(0)$ to be
$\mathfrak m_x$, so the stalk map factors through the residue field and the topological image and
sheaf map are compatible.  Without locality, the kernel could be a different prime and the stated
decomposition would fail.

If $X$ is a $k$-scheme and $K/k$ is fixed, the same proof works in the slice category.  One retains
only those $\kappa(x)\to K$ for which the structural maps from $k$ commute.  Thus
\[
X(K)\cong\coprod_{x\in X}
\operatorname{Hom}_{k\text{-compatible fields}}(\kappa(x),K).
\]
\end{solution}

% VI21-CHALLENGE-SOLUTION: C02
\paragraph{Solution to Challenge VI.21.C2.}
\begin{solution}
Define $y\preceq x$ when $y$ is a specialization of $x$, i.e.
$y\in\overline{\{x\}}$.  If $f:X\to Y$ is a scheme morphism and $y\preceq x$, continuity gives
\[
f(y)\in f(\overline{\{x\}})
\subseteq\overline{\{f(x)\}},
\]
so $f(y)\preceq f(x)$.  Thus the underlying point map is monotone for the specialization preorder.

Distinct comparable points may collapse whenever the morphism identifies their images.  The
simplest example is the structure morphism
\[
\Spec k[t]\longrightarrow\Spec k
\]
for a field $k$: the generic point $(0)$ and every closed point $(t-a)$ are distinct and comparable
through specialization, but all map to the unique point of $\Spec k$.  More generally, for an
affine morphism induced by $A\to B$, two primes $\mathfrak q\subsetneq\mathfrak q'$ of $B$ collapse
exactly when their contractions to $A$ coincide.  Monotonicity permits equality in the target, so
such collapse is fully compatible with functoriality.
\end{solution}

% VI21-CHALLENGE-SOLUTION: C03
\paragraph{Solution to Challenge VI.21.C3.}
\begin{solution}
Several examples separate topology from scheme structure.

First, $\Spec k$ and $\Spec k[\varepsilon]/(\varepsilon^2)$ are both one-point spaces, but the
second stalk contains a nonzero nilpotent.  Second,
\[
\Spec k[\varepsilon]/(\varepsilon^2)
\quad\text{and}\quad
\Spec k[\varepsilon]/(\varepsilon^3)
\]
again have the same one-point topology and the same residue field $k$, but their nilpotent
thickenings differ: the nilradical has square zero in the first ring and nonzero square in the
second.  Their global rings therefore cannot be isomorphic.

Third, $\Spec\mathbb R$ and $\Spec\mathbb C$ are homeomorphic one-point spaces with no nilpotents,
but their unique residue fields are different.  Hence even after reducedness is imposed, topology
does not recover the field carried by a point.

These examples show exactly what the underlying point space forgets: nilpotent infinitesimal
structure, local-ring multiplication, and residue-field data.  The structure sheaf restores this
information.
\end{solution}

% VI21-CHALLENGE-SOLUTION: C04
\paragraph{Solution to Challenge VI.21.C4.}
\begin{solution}
For every $x\in X$,
\[
\pi_X(x)=\ker(\mathbb Z\to\kappa(x)),
\]
so the underlying map $|X|\to|\Spec\mathbb Z|$ sends characteristic-zero points to the generic
point $(0)$ and characteristic-$p$ points to the closed point $(p)$.  Since every scheme morphism
preserves specialization, the structure morphism does as well.

The scheme $\Spec\mathbb Z$ itself already contains both kinds of points: its generic point $(0)$
has residue field $\mathbb Q$ of characteristic zero, while each closed point $(p)$ has residue
field $\mathbb F_p$.  Moreover $(p)$ is a specialization of $(0)$, matching the specialization
order in the target.  The same mixed-characteristic portrait appears in
$\Spec\mathbb Z[t]$: the generic point has characteristic zero, while primes such as $(p,t)$ have
characteristic $p$.

Thus the absolute base organizes a scheme into arithmetic strata at the level of points.  This
pointwise portrait is weaker than factorization of the whole scheme: for example
$\Spec(\mathbb Z/p^2\mathbb Z)$ has only characteristic-$p$ residue fields but does not factor
through $\Spec\mathbb F_p$.
\end{solution}

% VI21-CHALLENGE-SOLUTION: C05
\paragraph{Solution to Challenge VI.21.C5.}
\begin{solution}
Let $x\in X$ and choose an affine neighbourhood $U=\Spec A$ with
$x\leftrightarrow\mathfrak p$.  Under the identification
$\OO_{X,x}\cong A_{\mathfrak p}$, the canonical local-scheme morphism becomes
\[
\Spec A_{\mathfrak p}\longrightarrow\Spec A\hookrightarrow X.
\]
Its image inside $U$ is the set of contractions of primes of $A_{\mathfrak p}$, namely
\[
\{\mathfrak q\in\Spec A:\mathfrak q\subseteq\mathfrak p\}.
\]
These are exactly the generalizations of $x$ in $U$.

It remains to show that no generalization is missed globally.  If $y$ is a generalization of $x$,
then $x\in\overline{\{y\}}$.  Every open neighbourhood of $x$ must therefore contain $y$;
otherwise the closed complement would contain $y$ and hence its closure, contradicting
$x\in U$.  In particular $y\in U$.  Hence every global generalization is already visible in the
affine calculation above.

If a second affine chart $V$ contains $x$, then every generalization of $x$ lies in
$U\cap V$, and both chart descriptions identify the image with that same intrinsic set.  On a
smaller affine neighbourhood inside $U\cap V$, the localization maps agree through the canonical
stalk $\OO_{X,x}$.  Therefore the description is independent of the chart and
\[
\operatorname{im}(\Spec\OO_{X,x}\to X)
=\{\text{generalizations of }x\}.
\]
\end{solution}

